% Preamble integrasi untuk pembaca pilihan CRing berbahasa Indonesia.
% Berkas ini merupakan lapisan edisi; keenam fragmen terjemahan tetap terpisah.

\usepackage{fontspec}
\defaultfontfeatures{Ligatures=TeX}
\usepackage{amsmath}
\usepackage{amssymb}
\usepackage{amsfonts}
\usepackage{amscd}
\usepackage{mathrsfs}
\usepackage{hyperref}
\usepackage[amsmath,thmmarks]{ntheorem}
\usepackage[capitalize,noabbrev]{cleveref}
\usepackage{url}
\usepackage[top=1.15in,bottom=1.15in,left=1.15in,right=1.15in]{geometry}
\usepackage{fancyhdr}
\usepackage{verbatim}
\usepackage{enumerate}
\usepackage{makeidx}
\usepackage{longtable}
\usepackage{booktabs}
\usepackage{xcolor}
\usepackage{microtype}
\usepackage{titlesec}
\usepackage[all]{xy}

\makeindex
\renewcommand{\contentsname}{Daftar Isi}
\renewcommand{\bibname}{Daftar Pustaka}
\renewcommand{\indexname}{Indeks}
\renewcommand{\figurename}{Gambar}
\renewcommand{\tablename}{Tabel}
\setlength{\emergencystretch}{3em}
\hypersetup{
  colorlinks=true,
  linkcolor=blue!45!black,
  citecolor=green!35!black,
  urlcolor=blue!55!black,
  pdftitle={Pilihan CRing: Pembaca Bahasa Indonesia},
  pdfauthor={CRing Project; OpenAI Codex gpt-5.6-sol (Ultra)},
  pdfsubject={Enam pilihan aljabar komutatif dengan provenance GFDL yang eksplisit},
  pdfkeywords={aljabar komutatif, CRing, Bahasa Indonesia, GFDL}
}

\def\name{Pilihan CRing --- Bahasa Indonesia}
\pagestyle{fancy}
\fancyhf{}
\fancyhead[CE,CO]{\itshape \name}
\fancyfoot[CE,CO]{\thepage}
\renewcommand{\headrulewidth}{0pt}
\renewcommand{\footrulewidth}{0pt}
\setlength{\headheight}{14pt}

\titleformat{\chapter}[display]
  {\LARGE\bfseries\centering}{Bab \thechapter}{3pt}{}[\vspace{3pt}\hrule]
\titleformat{\section}{\Large\bfseries}{\S\,\thesection}{10pt}{}

% Makro matematika yang dipakai fragmen pilihan. Definisi ini mengikuti
% antarmuka sumber CRing, tetapi berada di preamble integrasi yang mandiri.
\newcommand{\lecture}[1]{}
\newcommand{\rad}{\operatorname{Rad}}
\newcommand{\im}{\operatorname{Im}}
\newcommand{\proj}{\operatorname{Proj}}
\renewcommand{\hom}{\operatorname{Hom}}
\newcommand{\id}{\operatorname{id}}
\providecommand{\cal}[1]{\mathcal{#1}}
\renewcommand{\cal}[1]{\mathcal{#1}}
\newcommand{\et}{\mathrm{Et}}
\newcommand{\supp}{\operatorname{supp}}
\renewcommand{\k}{\kappa}
\newcommand{\ann}{\operatorname{Ann}}
\newcommand{\ass}{\operatorname{Ass}}
\newcommand{\pic}{\operatorname{Pic}}
\newcommand{\cart}{\operatorname{Cart}}
\newcommand{\weil}{\operatorname{Weil}}
\newcommand{\krdim}{\operatorname{Krdim}}
\newcommand{\emdim}{\operatorname{Emdim}}
\newcommand{\Hom}{\operatorname{Hom}}
\renewcommand{\a}{{a}}
\renewcommand{\b}{{b}}
\renewcommand{\d}{{d}}
\newcommand{\e}{{e}}
\newcommand{\g}{{g}}
\newcommand{\het}{\operatorname{height}}
\newcommand{\f}{{f}}
\newcommand{\add}[1]{\textbf{AKAN DITAMBAHKAN:} #1}
\newcommand{\spec}{\operatorname{Spec}}
\newcommand{\cont}{\operatorname{Cont}}
\DeclareMathOperator{\Emb}{Emb}
\DeclareMathOperator{\Aut}{Aut}
\DeclareMathOperator{\Char}{char}
\DeclareMathOperator{\Gal}{Gal}
\DeclareMathOperator{\Tor}{Tor}
\DeclareMathOperator{\idem}{Idem}
\newcommand{\res}{\operatorname{Res}}
\newcommand{\coker}{\operatorname{Coker}}
\newcommand{\hearts}{\heartsuit}
\newcommand{\Ext}{\operatorname{Ext}}
\newcommand{\Tr}{\operatorname{Tr}}
\newcommand{\Specm}{\operatorname{Specm}}
\newcommand{\Ann}{\operatorname{Ann}}
\newcommand{\tor}{\operatorname{Tor}}
\newcommand{\ext}{\operatorname{Ext}}
\newcommand{\End}{\operatorname{End}}
\newcommand{\ob}{\operatorname{ob}}
\newcommand{\ad}{\mathrm{ad}}
\renewcommand{\log}{\mathrm{log}}
\renewcommand{\dim}{\mathrm{dim}}
\newcommand{\dd}[2]{\frac{\partial #1}{\partial #2}}
\newcommand{\gr}{\operatorname{gr}}
\newcommand{\st}{\ss:\ss}
\newcommand{\Sym}{\mathrm{Sym}}
\newcommand{\depth}{\operatorname{depth}}

\newenvironment{xyxy}[1]
  {\vspace{10pt}\begin{center}\xymatrix{#1}\end{center}\vspace{12pt}}
  {}
\newenvironment{qu}[2]
  {\begin{list}{}{\setlength\leftmargin{#1}\setlength\rightmargin{#2}}\item[]\footnotesize}
  {\end{list}}
\newenvironment{random}[1]{\par\vspace{10pt}\noindent\textbf{#1}}{}
\newenvironment{liszt}[1]
  {\begin{list}{\labelitemi}{\setlength\leftmargin{#1em}\setlength\itemsep{0.1in}}}
  {\end{list}}
\newcommand{\sqxx}[4]{\begin{xyxy}{#1 \ar[r] \ar[d] & #2 \ar[d] \\ #3 \ar[r] & #4}\end{xyxy}}

% Lingkungan teorema memakai nama Indonesia tanpa mengubah nama lingkungan
% yang terdapat dalam keenam fragmen.
\newtheorem{theorem}{Teorema}[section]
\newtheorem{lemma}[theorem]{Lema}
\newtheorem{sublemma}[theorem]{Sublema}
\newtheorem{corollary}[theorem]{Korolari}
\newtheorem{proposition}[theorem]{Proposisi}
\theorembodyfont{\normalfont}
\newtheorem{definition}[theorem]{Definisi}
\newtheorem*{remark}{Catatan}
\newtheorem{example}[theorem]{Contoh}
\newtheorem{exercise}{Latihan}[chapter]
\newtheorem*{solution}{Penyelesaian}
\newtheorem*{question}{Pertanyaan}
\newtheorem*{problem}{Masalah}
\newtheorem*{dbend}{Peringatan}
\newtheorem*{prediction}{Prediksi}
\theoremheaderfont{\itshape}
\theoremseparator{.}
\theoremsymbol{\ensuremath{\blacktriangle}}
\newtheorem*{proof}{Bukti}

\crefname{theorem}{teorema}{teorema}
\crefname{lemma}{lema}{lema}
\crefname{corollary}{korolari}{korolari}
\crefname{proposition}{proposisi}{proposisi}
\crefname{definition}{definisi}{definisi}
\crefname{example}{contoh}{contoh}
\crefname{exercise}{latihan}{latihan}
\crefname{section}{bagian}{bagian}
\crefname{subsection}{subbagian}{subbagian}

% Referensi berikut memang berada di luar enam rentang terpilih. Registri ini
% mengubahnya menjadi tautan yang jujur ke daftar dependensi sumber lengkap,
% bukan tanda tanya LaTeX. Pemetaan persisnya juga dicatat dalam manifest JSON.
\makeatletter
\let\CRingNativeCref\cref
\newcommand{\CRingDeclareExternalReference}[2]{%
  \expandafter\gdef\csname CRingExternalRef@#1\endcsname{#2}}
\CRingDeclareExternalReference{Jacobson}{latihan tentang radikal Jacobson dalam sumber lengkap}
\CRingDeclareExternalReference{def-radical-ideal}{definisi ideal radikal dalam sumber lengkap}
\CRingDeclareExternalReference{anycontainedinmaximal}{proposisi ideal maksimal dalam sumber lengkap}
\CRingDeclareExternalReference{primeavoidance}{teorema penghindaran prima dalam sumber lengkap}
\CRingDeclareExternalReference{assmfinite}{proposisi kefinitan prima terkait dalam sumber lengkap (label sumber: \texttt{finiteassm})}
\CRingDeclareExternalReference{spec}{bab tentang spektrum dalam sumber lengkap}
\CRingDeclareExternalReference{nilpcriterion}{kriteria nilpotensi dalam sumber lengkap}
\CRingDeclareExternalReference{noetherian}{bab tentang gelanggang Noether dalam sumber lengkap}
\CRingDeclareExternalReference{finitelygeneratedintegral}{proposisi tentang aljabar integral terbangkitkan hingga dalam sumber lengkap}
\CRingDeclareExternalReference{subsectiondimension}{subbagian teori dimensi ruang topologis dalam sumber lengkap}
\newcommand{\CRingResolvedReference}[1]{%
  \@ifundefined{CRingExternalRef@#1}
    {\CRingNativeCref{#1}}
    {\hyperref[cring-external-reference-registry]{\csname CRingExternalRef@#1\endcsname}}}
\renewcommand{\cref}[1]{\CRingResolvedReference{#1}}
\newcommand{\rref}[1]{\CRingResolvedReference{#1}}
\makeatother
